% Глава: 19-docker-basics
% Статус: наполнено (первый прогон).
\chapter{Docker: основы}
\label{ch:19-docker-basics}

\section{Контейнер и виртуальная машина}

\textbf{Виртуальная машина} виртуализует
\emph{оборудование}: внутри гостевая ОС с ядром,
драйверами, собственным пространством.
\textbf{Контейнер} изолирует \emph{процессы}: все
контейнеры на хосте делят одно ядро Linux, но имеют
отдельные файловые системы, сетевые
пространства и PID-пространства.

На курсе Docker-демон (хост) --- это ваша ВМ на сервере
\fpath{pglab.nsu.ru}: контейнеры изолируются внутри вашей системы,
разделяя её ядро и диск.

\begin{figure}[htb]
  \centering
  \begin{tikzpicture}[node distance=0.6cm]
    \node[layer, minimum width=6cm, minimum height=2.6cm] (vmhost) {};
    \node[font=\small, above=2pt of vmhost] {Виртуальные машины};
    \node[process, minimum width=2.4cm] (vm1) at (vmhost) {}; 
    \node[font=\small] at (vm1) {ВМ: гостевая ОС + ядро};
    \node[layer, below=1.6cm of vmhost, minimum width=6cm, minimum height=2.6cm] (cghost) {};
    \node[font=\small, above=2pt of cghost] {Контейнеры};
    \node[ext, minimum width=1.5cm] (c1) at ([yshift=-0.5cm]cghost.west) [xshift=0.7cm] {app};
    \node[ext, right=0.3cm of c1, minimum width=1.5cm] (c2) {app};
    \node[ext, right=0.3cm of c2, minimum width=1.5cm] (c3) {app};
    \node[font=\small, below=3pt of c2] {общее ядро хоста};
  \end{tikzpicture}
  \caption{ВМ изолирует целую ОС, контейнеры --- процессы поверх общего ядра.}
  \label{fig:vm-vs-ct}
\end{figure}

Преимущества контейнеров: быстрый запуск (секунды),
малый размер (сотни МБ), переносимость «от
разработчика до продакшена».

\section{Архитектура Docker}

\begin{itemize}
  \item \textbf{Docker Engine} (демон \cmd{dockerd}) ---
    управляющий процесс на хосте;
  \item \textbf{клиент} (\cmd{docker}) --- команда,
    которой вы пользуетесь; общается с демоном по
    API;
  \item \textbf{образ} (image) --- неизменяемый набор
    файлов + метаданные; \textbf{контейнер} ---
    работающий экземпляр образа.
\end{itemize}

\section{Ключевые команды}

\begin{minted}[bgcolor=gray!6]{bash}
docker run hello-world        # скачать и запустить тестовый образ
docker run -it ubuntu bash    # запустить интерактивный контейнер
docker ps                     # работающие контейнеры
docker ps -a                  # все контейнеры, включая остановленные
docker images                 # локальные образы
docker exec -it <id> bash     # выполнить команду в работающем контейнере
docker logs <id>              # stdout/stderr контейнера
docker stop <id>              # остановить контейнер
docker start <id>             # запустить остановленный контейнер
docker rm <id>                # удалить остановленный контейнер
docker rmi <image>            # удалить образ
\end{minted}

\cmd{docker run} с флагами:
\begin{minted}[bgcolor=gray!6]{bash}
docker run -d -p 8080:80 --name web nginx   # в фоне, маппинг портов
docker run --rm -it alpine sh               # удалить при выходе
docker run --name app myimage:1.0           # по имени и тегу образа
\end{minted}

\section{Образы и реестры}

\textbf{Реестр} (registry) --- хранилище образов;
основной публичный --- Docker Hub.
\textbf{Тег} (tag) --- метка версии в имени образа:
\cmd{nginx:1.25}, \cmd{nginx:latest}.
\begin{minted}[bgcolor=gray!6]{bash}
docker pull nginx:1.25           # скачать конкретную версию
docker tag myimage:1.0 me/myimage:1.0   # переименовать для push
docker push me/myimage:1.0       # загрузить, после login
\end{minted}

\section{Типичные ошибки}
\begin{itemize}
  \item \emph{permission denied} на \fpath{/var/run/docker.sock}
    --- пользователь не в группе \cmd{docker}; добавьте
    (\cmd{sudo usermod -aG docker \$USER}), перезайдите.
  \item Контейнер сразу завершился --- проверьте
    \cmd{docker logs}; часто причина --- отсутствующий
    \cmd{CMD} в Dockerfile.
  \item Порт «не отвечает» --- не задан \cmd{-p host:container}.
  \item Изменения в контейнере потерялись после
    перезапуска --- данные хранятся в слое контейнера;
    для сохранения нужны volumes (\nameref{ch:21-docker-data}).
\end{itemize}

\section{Связь с остальным курсом}

Создание образа --- Dockerfile
(\nameref{ch:20-dockerfile}); хранение данных ---
\nameref{ch:21-docker-data}; несколько контейнеров ---
\nameref{ch:22-docker-compose}; ограничения ресурсов ---
\nameref{ch:23-docker-resources}. Внутрь контейнера
можно войти \cmd{docker exec} и поработать
в tmux (\nameref{ch:10-tmux}) или Vim
(\nameref{ch:09-vim}).

\section{Задачи для самопроверки}
\begin{enumerate}
  \item Чем контейнер отличается от виртуальной
    машины?
  \item Что такое образ и что такое контейнер?
  \item Как запустить интерактивный контейнер с
    bash и как из него выйти?
  \item Как посмотреть логи завершившегося
    контейнера?
  \item Зачем нужны теги образов?
\end{enumerate}
